\subsection{Dalil Mensyukuri Nikmat}
\begin{enumerate}
\item \textit{Wa iż ta'ażżana rabbukum la'in syakartum la'azīdannakum wa la'in kafartum inna ‘ażābī lasyadīd.} \\
  \textbf{Artinya: }(Ingatlah) ketika tuhanmu memaklumkan, sesungguhnya jika kamu bersyukur, niscaya aku akan menambah (nikmat) kepadmu, tetapi jika kamu mengingkari (nikmat-ku), sesungguhnya azab-ku benar-benar sangat keras. \textbf{Ibrahim (7)}
  
\item \textit{Wallāhu akhrajakum mim buṭūni ummahātikum lā ta‘lamūna syai'aw-wa ja‘ala lakumus-sam‘a wal-abṣāra wal-af'idah, la‘allakum tasykurūn.} \\
  \textbf{Artinya: }Allah mengeluarkan kamu dari perut ibumu dalam keadaan tidak mengetahui sesuatu pun dan dia menjadikan bagi kamu pendengaran, penglihatan, dan hati nurani agar kamu bersyukur. \textbf{An-Nahl (78)}

\item \textit{‘An Abī Hurairata raḍiyallāhu ‘anhu qāla: qāla Rasūlullāhi ṣallallāhu ‘alaihi wa sallam: Man lā yasykurin-nāsa lā yasykurillāh.} \\
  \textbf{Artinya: }Dari Abu Hurairah R.A ia berkata: Rasulullah SAW bersabda: ``Barangsiapa yang tidak bersyukur (berterima kasih) kepada manusia, maka ia tidak bersyukur kepada Allah.'' \textbf{(HR Tirmidzi \& Ahmad)}
\end{enumerate}


Secara istilah syukur adalah sikap mengakui dan menyadari bahwa segala nikmat yang diterima berasal dari Allah SWT, kemudian menggunakan nikmat tersebut sesuai dengan kehendaknya.


\subsubsection{Unsur Utama Syukur}
\begin{enumerate}
\item \textbf{Syukur dengan hati: }Yaitu menyadari dan meyakini semua nikmat berasal dari Allah SWT.
\item \textbf{Syukur dengan lisan: }Yaitu mengucapkan pujian kepada allah dan menggunakan perkataan yang baik.
\item \textbf{Syukur dengan perbuatan: }Yaitu menggunakan nikmat Allah untuk berbuat kebaikan.
\end{enumerate}


\subsubsection{Cara Mensyukuri Nikmat}
\begin{enumerate}
\item Bersyukur atas nikmat kesehatan dengan cara menjaga kebersihan, makan makanan yang baik, berolahraga, rajin beribadah, tidak merusak tubuh.
\item Bersyukur atas nikmat ilmu dengan cara rajin belajar menggunakan ilmu untuk kebaikan, menghormati guru, sharing ilmu pengetahuan.
\item Bersyukur atas nikmat harta dengan cara:
  \begin{enumerate}
  \item Mencari nafkah dengan cara yang halal.
  \item Tidak boros.
  \item Bersedekah.
  \item Membantu orang yang membutuhkan.
  \end{enumerate}
\end{enumerate}

